\chapter{Software Architecture}
\label{chap:architecture}

\section{Architecture Overview}
\label{sec:architecture-overview}\index{architecture}\index{login}\index{guest}\index{Keycloak}\index{Caddy}

A question crosses eight layers: from the web chat to access control, the input guardrail, the
LangGraph\index{LangGraph} orchestrator~\cite{langgraph}, the sources, the generation and the
output guardrail, up to the answer or an operator. Login is optional: a guest gets general answers
from the public documents, a customer also gets answers on their own policy. Everything is free:
one server, plus Groq's free \acs{api} for the two large models (Figure~\ref{fig:overview-strip}).

% Figure 2.1: a question at a glance, from the chat to the answer (details in the next figure)
\begin{figure}[H]
  \centering
  \begin{tikzpicture}[font=\scriptsize,
      lay/.style={rounded corners=3pt, text width=1.5cm, minimum height=1.55cm, align=center,
                  inner sep=2pt},
      flow/.style={-{Stealth[length=1.8mm]}, accent, line width=0.7pt}]
    \newcommand\lgx[1]{\includegraphics[height=0.5cm, width=0.62cm, keepaspectratio]{logos/#1}}
    % \lay{index}{category}{icon or logo}{label}
    \newcommand\lay[4]{\node[lay, cat #2] (l#1) at ({(#1-1)*1.98},0)
      {#3\\[3pt]{\bfseries #4}};}
    \lay{1}{infra}{{\color{accent}\Large\faComments}}{Web chat}
    \lay{2}{sec}{\lgx{keycloak}}{Optional login}
    \lay{3}{sec}{\lgx{qwen}}{Input guardrail}
    \lay{4}{infra}{\lgx{langgraph}}{Orche\-strator}
    \lay{5}{data}{\lgx{opensearch}\,\lgx{mcp}}{Sources}
    \lay{6}{model}{\lgx{openai}}{Generation}
    \lay{7}{sec}{{\color{catsec}\Large\faCheckDouble}}{Output guardrail}
    \lay{8}{infra}{\lgx{chatwoot}}{Answer or operator}
    \foreach \a/\b in {1/2,2/3,3/4,4/5,5/6,6/7,7/8}
      \draw[flow] (l\a) -- (l\b);
  \end{tikzpicture}
  \caption{A question at a glance.}
  \label{fig:overview-strip}
\end{figure}


\section{Components}
\label{sec:components}\index{policy server}\index{token exchange}\index{core insurance system}

Most components in Figure~\ref{fig:architecture} are standard pieces chosen for being free and
well documented; the one written specifically for this project is the policy server. It is a
server based on the Model Context Protocol\index{Model Context Protocol}~\cite{mcp} that stands in
front of the company's core insurance system and exposes only four read-only operations:

\begin{itemize}
  \item the list of the policies of the logged-in customer, with product, edition of the
    conditions, status and dates, and the insured object masked;
  \item the guarantees, limits, deductibles and premium of one policy;
  \item the instalments of a policy, with their due dates and status;
  \item the open claims of a policy, with their stage and the contact person.
\end{itemize}

Each operation protects the data in its own way. A customer only ever sees policies that belong to
the subject of the login token: asking for someone else's policy returns \enquote{not found}, rather
than \enquote{forbidden}, so as not to confirm that the policy exists. Bank details are cut to their
last four digits, and the claims never include reserves or internal notes. The token used by the
policy server is not the one of the chat: Keycloak issues a new one, limited to this purpose,
through the standard token exchange~\cite{keycloak2026tokenexchange}.

% Figure 2.1: the eight layers crossed by a question, with the logo of every technology.
% Tools sit in three columns, at x = 4.85, 9.25 and 13.65.
\begin{figure}[p]
  \centering
  \begin{tikzpicture}[
      font=\scriptsize,
      layer/.style={draw=black!30, fill=black!4, rounded corners=4pt},
      band/.style={draw=black!45, dashed, rounded corners=4pt},
      tool/.style={rounded corners=2pt, minimum height=0.85cm, inner sep=2pt},
      layername/.style={anchor=west, text width=2.3cm, align=left, font=\footnotesize\bfseries},
      grouptitle/.style={font=\scriptsize\bfseries, text=black!65},
      hop/.style={-{Stealth[length=2mm]}, black!55, line width=0.6pt},
      hoplabel/.style={anchor=west, font=\scriptsize\itshape, text=black!70},
      swatch/.style={rounded corners=1.5pt, minimum width=0.55cm, minimum height=0.32cm},
    ]
    % \tool{category}{x}{y}{logo or icon}{name}{role}
    \newcommand\tool[6]{\node[tool, cat #1] at (#2,#3) {%
      \begin{tabular}{@{}>{\centering\arraybackslash}m{0.55cm}@{\hspace{3pt}}
                         >{\raggedright\arraybackslash}m{3.45cm}@{}}
        #4 & \textbf{#5}\newline #6
      \end{tabular}};}
    \newcommand\lgo[1]{\logo[0.5cm]{#1}}
    \newcommand\ico[1]{{\color{accent}\Large #1}}
    % \layer{top}{bottom}{number and name}
    \newcommand\layer[3]{%
      \draw[layer] (0,#1) rectangle (15.85,#2);
      \node[layername] at (0.15,{(#1+#2)/2}) {#3};}
    % \hop{from y}{to y}{what travels between two layers}
    \newcommand\hop[3]{%
      \draw[hop] (0.9,#1) -- (0.9,#2);
      \node[hoplabel] at (1.05,{(#1+#2)/2}) {#3};}

    \layer{0.6}{-0.6}{1\enspace Channels}
    \tool{infra}{4.85}{0}{\ico{\faComments}}{Web chat}{HTML + JS, served by FastAPI}
    \hop{-0.6}{-1.1}{message, + JWT after login}

    \layer{-1.1}{-2.3}{2\enspace Access}
    \tool{infra}{4.85}{-1.7}{\lgo{caddy}}{Caddy}{HTTPS, reverse proxy}
    \tool{sec}{9.25}{-1.7}{\lgo{keycloak}}{Keycloak}{optional login, token exchange}
    \hop{-2.3}{-2.8}{verified identity}

    \layer{-2.8}{-4.0}{3\enspace Input guardrail}
    \tool{model}{4.85}{-3.4}{\lgo{qwen}}{Qwen3Guard-Gen-0.6B}{jailbreaks, banned content}
    \tool{sec}{9.25}{-3.4}{\lgo{presidio}}{Presidio}{masks personal data}
    \hop{-4.0}{-4.5}{admitted question}

    \layer{-4.5}{-5.7}{4\enspace Orchestrator}
    \tool{infra}{4.85}{-5.1}{\lgo{langgraph}}{LangGraph}{graph of 13 nodes}
    \tool{infra}{9.25}{-5.1}{\lgo{fastapi}}{FastAPI + SSE}{API, streamed states}
    \tool{model}{13.65}{-5.1}{\lgo{openai}}{gpt-oss-20b}{intent, rewording \textbullet\ Groq}
    \hop{-5.7}{-6.2}{query, filters, service token}

    \layer{-6.2}{-9.6}{5\enspace Sources\\[2pt]{\mdseries\color{black!60}read in parallel}}
    \node[grouptitle] at (4.85,-6.45) {Knowledge};
    \node[grouptitle] at (9.25,-6.45) {Customer data};
    \node[grouptitle] at (13.65,-6.45) {Memory};
    \tool{data}{4.85}{-7.125}{\lgo{opensearch}}{OpenSearch}{BM25 + vectors, filters}
    \tool{model}{4.85}{-8.075}{\lgo{qwen}}{Qwen3-Embedding-0.6B}{clause vectors}
    \tool{model}{4.85}{-9.025}{\lgo{qwen}}{Qwen3-Reranker-0.6B}{reorders the top 40}
    \tool{data}{9.25}{-7.125}{\lgo{mcp}}{MCP policy server}{after login, read only}
    \tool{infra}{9.25}{-8.075}{\ico{\faBuilding}}{Core insurance system}{existing company APIs}
    \tool{data}{13.65}{-7.125}{\lgo{postgresql}}{PostgreSQL}{checkpoints, audit}
    \tool{data}{13.65}{-8.075}{\lgo{valkey}}{Valkey}{cache, Chatwoot queues}
    \hop{-9.6}{-10.1}{clauses + policy data}

    \layer{-10.1}{-11.3}{6\enspace Generation}
    \tool{model}{4.85}{-10.7}{\lgo{openai}}{gpt-oss-120b}{117B MoE \textbullet\ Groq}
    \hop{-11.3}{-11.8}{draft with citations}

    \layer{-11.8}{-13.0}{7\enspace Output guardrail}
    \tool{sec}{4.85}{-12.4}{\ico{\faCheckDouble}}{Validator}{citations, amounts, dates}
    \tool{model}{9.25}{-12.4}{\lgo{openai}}{gpt-oss-20b}{groundedness, criteria}
    \hop{-13.0}{-13.5}{verified answer}

    \layer{-13.5}{-14.7}{8\enspace Output}
    \tool{infra}{4.85}{-14.1}{\ico{\faIcon{file-alt}}}{Answer with sources}{citations, intent, confidence}
    \tool{infra}{9.25}{-14.1}{\lgo{chatwoot}}{Chatwoot}{operator console, corrections}

    % Services used by every layer
    \draw[band] (0,-15.3) rectangle (15.85,-18.7);
    \node[layername] at (0.15,-17.0)
      {Cross-cutting services\\[2pt]{\mdseries\color{black!60}used by every layer}};
    \node[grouptitle] at (4.85,-15.55) {Observability};
    \node[grouptitle] at (9.25,-15.55) {Quality};
    \node[grouptitle] at (13.65,-15.55) {Hosting};
    \tool{infra}{4.85}{-16.225}{\ico{\faStream}}{JSON logs}{a trace per node, masked}
    \tool{sec}{4.85}{-17.175}{\ico{\faLink}}{Audit chain}{hash-linked rows}
    \tool{infra}{9.25}{-16.225}{\lgo{ragas}}{Ragas}{faithfulness, relevancy}
    \tool{sec}{9.25}{-17.175}{\lgo{promptfoo}}{promptfoo}{red teaming, injection}
    \tool{infra}{9.25}{-18.125}{\lgo{pytest}}{pytest}{routes, scenarios}
    \tool{infra}{13.65}{-16.225}{\lgo{docker}}{Docker Compose}{every service, one file}
    \tool{infra}{13.65}{-17.175}{\lgo{oracle}}{Oracle Cloud}{free server, 4 cores, 24 GB}

    % Legend
    \node[swatch, cat model] at (0.4,-19.35) {};
    \node[anchor=west] at (0.75,-19.35) {open-weight model};
    \node[swatch, cat sec] at (3.8,-19.35) {};
    \node[anchor=west] at (4.15,-19.35) {security and controls};
    \node[swatch, cat data] at (7.7,-19.35) {};
    \node[anchor=west] at (8.05,-19.35) {data and knowledge};
    \node[swatch, cat infra] at (11.4,-19.35) {};
    \node[anchor=west] at (11.75,-19.35) {infrastructure and channels};
  \end{tikzpicture}
  \caption{The layers of the assistant, crossed by a question from top to bottom. Everything runs
    on one free server, except the two gpt-oss models (Groq's free \acs{api}).}
  \label{fig:architecture}
\end{figure}

\clearpage

\section{Component Interactions}
\label{sec:interactions}

The intent recognised by \code{classify} decides the route (Figure~\ref{fig:intents}). A turn is
one run of the graph (Figure~\ref{fig:graph}): after each node, a function written in code picks
the next one.\index{intent}

% Figure 2.3: the eight intents recognised by classify and the route each one takes
\begin{figure}[H]
  \centering
  \begin{tikzpicture}[
      font=\scriptsize,
      chip/.style={rounded corners=6pt, draw=catmodel!60, fill=catmodel!8, inner sep=3pt,
                   align=left, text width=2.75cm, anchor=west},
      dest/.style={rounded corners=3pt, draw=black!40, fill=white, inner sep=4pt, align=left,
                   text width=4.9cm},
      route/.style={-{Stealth[length=2mm]}, black!55, line width=0.6pt},
    ]
    \node[rounded corners=3pt, cat model, align=center, text width=1.7cm, minimum height=1.2cm]
      (cl) at (0.95,-2.5) {\logo[0.5cm]{openai}\\[2pt]{\ttfamily classify}};
    % \route{id}{y}{label}{meaning}{second label}{second meaning}{destination}
    \newcommand\intent[7]{%
      \node[chip] (a#1) at (2.5,#2) {{\ttfamily\bfseries #3}\\{\color{black!65}#4}};
      \ifx\relax#5\relax
        \coordinate (b#1) at (a#1.east);
      \else
        \node[chip] (c#1) at (5.75,#2) {{\ttfamily\bfseries #5}\\{\color{black!65}#6}};
        \coordinate (b#1) at (c#1.east);
      \fi
      \node[dest] (d#1) at (13.25,#2) {#7};
      \draw[route] (cl.east) -- (a#1.west);
      \draw[route] (b#1) -- (d#1.west);}
    \intent{1}{0}{polizza}{conditions of one's policy}{}{}
      {after login: \code{resolve_policy}, then \code{fetch_policy_data} and \code{retrieve} in
       parallel. A guest is asked to log in}
    \intent{2}{-1.3}{sinistro}{status of a claim}{pagamenti}{instalments, payments}
      {after login: \code{resolve_policy}, then \code{fetch_policy_data} with the claims or the
       instalments. A guest is asked to log in}
    \intent{3}{-2.6}{prodotto}{a product in general}{}{}
      {\code{retrieve} on the public documents, for guests and customers alike}
    \intent{4}{-3.8}{operatore}{asks for a person}{reclamo}{a complaint}
      {\code{handoff}; a complaint goes with high priority}
    \intent{5}{-5.0}{fuori\_ambito}{off-topic request}{saluto}{greeting}
      {\code{direct_reply}, a fixed text}
  \end{tikzpicture}
  \caption{The intents recognised by \texttt{classify} (labels as in the code) and their routes.}
  \label{fig:intents}
\end{figure}

% Figure 2.2: the LangGraph graph of a conversation turn; badges show the technology of each node
\begin{figure}[p]
  \centering
  \begin{tikzpicture}[
      font=\scriptsize,
      gnode/.style={rounded corners=3pt, text width=2.6cm, minimum height=0.8cm, align=center,
                    inner sep=2pt},
      term/.style={dashed},
      pill/.style={rounded corners=7pt, draw=black!55, fill=black!5, minimum width=1.4cm,
                   minimum height=0.45cm, font=\scriptsize\bfseries},
      route/.style={-{Stealth[length=2mm]}, black!60, line width=0.6pt},
      loop/.style={route, dashed},
      elabel/.style={font=\scriptsize, text=black!75, align=center, inner sep=1.5pt},
      badge/.style={circle, fill=white, draw=black!25, inner sep=1pt, minimum size=0.5cm},
    ]
    % \gnode{name}{category, extra styles}{x}{y}{node}{what it uses}
    \newcommand\gnode[6]{%
      \node[gnode, cat #2] (#1) at (#3,#4) {{\ttfamily #5}\\{\color{black!70}#6}};}
    % \badge{node}{logo}: the logo on the top-left corner of a node
    \newcommand\badge[2]{%
      \node[badge] at ([xshift=-1.6mm, yshift=1.6mm]#1.north west) {\includegraphics[height=0.3cm, width=0.34cm,
        keepaspectratio]{logos/#2}};}

    \node[pill] (start) at (4.6,0) {START};
    \gnode{ig}{sec}{4.6}{-1.2}{input\_guard}{Qwen3Guard 0.6B}
    \gnode{cl}{model}{4.6}{-2.75}{classify}{gpt-oss-20b}
    \gnode{dr}{infra, term}{0}{-2.75}{direct\_reply}{fixed replies}
    \gnode{rp}{data}{4.6}{-4.3}{resolve\_policy}{MCP, policy list}
    \gnode{ac}{infra, term}{0}{-4.3}{ask\_clarification}{which policy?}
    \gnode{fp}{data}{4.6}{-5.85}{fetch\_policy\_data}{MCP, read only}
    \gnode{rt}{data}{9.0}{-5.85}{retrieve}{search + rerank}
    \gnode{gr}{data}{4.6}{-7.4}{grade}{rerank threshold}
    \gnode{rq}{model}{9.0}{-7.4}{rewrite\_query}{gpt-oss-20b, glossary}
    \gnode{ge}{model}{4.6}{-8.95}{generate}{gpt-oss-120b}
    \gnode{ho}{infra, term}{12.9}{-8.95}{handoff}{Chatwoot + pause}
    \gnode{og}{sec}{4.6}{-10.5}{output\_guard}{validator + 20b}
    \gnode{fi}{infra}{4.6}{-12.05}{finalize}{sources, audit}
    \node[pill] (end) at (4.6,-13.3) {END};
    \badge{ig}{qwen} \badge{cl}{openai} \badge{rq}{openai}
    \badge{rp}{mcp} \badge{fp}{mcp} \badge{rt}{opensearch}
    \badge{ge}{openai} \badge{og}{openai} \badge{ho}{chatwoot}
    \badge{fi}{postgresql}

    \draw[route] (start) -- (ig);
    \draw[route] (ig) -- node[elabel, right] {admitted} (cl);
    \draw[route] (ig.west) -| node[elabel, above, pos=0.25] {blocked} (dr.north);
    \draw[route] (cl.west) -- node[elabel, above] {greeting,\\off-topic,\\no login} (dr.east);
    \draw[route] ([yshift=2mm]cl.east) -| (ho.north);
    \node[elabel, above] at (11.0,-2.55) {operator, complaint};
    \draw[route] ([yshift=-2mm]cl.east) -| (rt.north);
    \node[elabel, below] at (7.6,-2.95) {product};
    \draw[route] (cl) -- node[elabel, right] {policy, claim,\\payment} (rp);
    \draw[route] (rp.west) -- node[elabel, above] {no single\\policy} (ac.east);
    \draw[route] (rp) -- node[elabel, right] {policy found} (fp);
    \draw[route] (rp.east) -- node[elabel, above] {in parallel} (9.0,-4.3);
    \fill[black!60] (9.0,-4.3) circle (1.6pt);
    \draw[route] (fp) -- (gr);
    \draw[route] (fp.west) -- ++(-0.6,0) |- ([yshift=1.5mm]ge.west);
    \node[elabel, left] at (2.6,-7.4) {claim,\\payment};
    \draw[route] ([xshift=-6mm]rt.south) -- ++(0,-0.2) -| ([xshift=6mm]gr.north);
    \draw[route] ([yshift=1.5mm]gr.east) -- node[elabel, above, pos=0.35] {context\\too weak}
      ([yshift=1.5mm]rq.west);
    \draw[loop] ([xshift=6mm]rq.north) -- node[elabel, right] {max 2} ([xshift=6mm]rt.south);
    \draw[route] (gr) -- node[elabel, right] {enough\\context} (ge);
    \draw[route] ([yshift=-1.5mm]gr.east) -- ++(0.45,0) |- (ho.west);
    \node[elabel, above] at (9.4,-8.95) {no relevant clause};
    \draw[route] (ge) -- (og);
    \draw[loop] (og.west) -- ++(-0.3,0) |- ([yshift=-1.5mm]ge.west);
    \node[elabel, left] at (2.85,-9.8) {regenerate,\\max 1};
    \draw[route] (og) -- node[elabel, right] {verified} (fi);
    \draw[route] (og.east) -| (ho.south);
    \node[elabel, above] at (9.4,-10.5) {not verifiable};
    \draw[route] (fi) -- (end);

    % Legend: what colours, borders, arrows and badges mean
    \begin{scope}[yshift=-14.3cm]
      \node[rounded corners=4pt, draw=black!20, fill=black!3, minimum width=15.4cm,
            minimum height=1.55cm] at (6.45,-0.45) {};
      \newcommand\sw[3]{\node[rounded corners=2pt, cat #1, minimum width=0.6cm,
        minimum height=0.32cm] at (#2,0) {}; \node[anchor=west] at (#2+0.4,0) {#3};}
      \sw{model}{-0.6}{language model}
      \sw{sec}{3.3}{check}
      \sw{data}{5.6}{data and search}
      \sw{infra}{9.0}{output}
      \node[rounded corners=2pt, cat infra, dashed, minimum width=0.6cm, minimum height=0.32cm]
        at (11.3,0) {};
      \node[anchor=west] at (11.7,0) {ends the turn};
      \draw[route] (-0.9,-0.85) -- (-0.1,-0.85);
      \node[anchor=west] at (0,-0.85) {route chosen by code};
      \draw[loop] (3.0,-0.85) -- (3.8,-0.85);
      \node[anchor=west] at (3.9,-0.85) {loop with a maximum};
      \fill[black!60] (7.6,-0.85) circle (1.6pt);
      \node[anchor=west] at (7.8,-0.85) {nodes in parallel};
      \node[badge] at (10.9,-0.85) {\includegraphics[height=0.3cm, width=0.34cm,
        keepaspectratio]{logos/openai}};
      \node[anchor=west] at (11.25,-0.85) {technology used};
    \end{scope}
  \end{tikzpicture}
  \caption{The graph of a conversation turn.}
  \label{fig:graph}
\end{figure}


\subsection{Routing Rules}
\label{sec:routing}\index{confidence}\index{policy selection}\index{rewording}

Three decisions of the graph are plain rules, not model outputs, and their counters restart at
every turn.

\paragraph{Which policy?} A customer may hold several policies. The graph uses the one the customer
names, by its number or by the plate of the vehicle; otherwise, the only active one; otherwise, the
one discussed in the previous turn. If none of these applies, it asks the customer.

\paragraph{Is the context enough?} The reranker gives each retrieved clause a score. Clauses below
a minimum threshold are dropped, the others are kept, and those above a higher threshold are
considered strong. The confidence shown with the answer follows from them: it is high when at least
one clause is strong and medium when at least two are kept. Otherwise the context is too weak to
answer. The thresholds are calibrated on the golden set of test questions.

\paragraph{How many attempts?} When the context is too weak, the question is reworded with
insurance synonyms, such as \textit{furto} and \textit{sottrazione} (theft), \textit{franchigia} and
\textit{scoperto} (deductible), \textit{sinistro} and \textit{evento} (claim), and searched again, at
most twice. When a draft fails the output check, it is regenerated once. If the problem persists,
the conversation goes to an operator.

\needspace{7cm}
\subsection{Operator Corrections}
\label{sec:corrections}\index{operator corrections}\index{unanswered questions}

What the operators correct, and what the assistant could not answer, comes back as better
answers, through a human review (Figure~\ref{fig:feedback}).

% Figure 2.9: how the operators' corrections improve the assistant
\begin{figure}[H]
  \centering
  \begin{tikzpicture}[font=\scriptsize,
      fbstep/.style={rounded corners=3pt, text width=2.9cm, minimum height=1.45cm, align=left,
                   inner sep=4pt},
      fbout/.style={rounded corners=3pt, draw=good!60, fill=good!7, text width=3.0cm,
                  minimum height=0.95cm, align=left, inner sep=3pt},
      flow/.style={-{Stealth[length=2.2mm]}, accent, line width=0.8pt}]
    % \fb{name}{x}{category}{icon}{title}{text}
    \newcommand\fb[6]{\node[fbstep, cat #3] (#1) at (#2,0)
      {{\footnotesize\bfseries #4\enspace #5}\\[2pt]#6};}
    \fb{a}{1.6}{infra}{{\color{accent}\faComments}}{Answer}
      {with its sources, the intent and the confidence}
    \fb{o}{5.55}{infra}{\logo{chatwoot}}{Operator}
      {in Chatwoot, marks a wrong answer and writes the right one, with the reason}
    \fb{l}{9.5}{data}{\logo{postgresql}}{Log}
      {corrections and questions left without an answer, masked}
    \fb{r}{13.45}{sec}{{\color{catsec}\faUsers}}{Review}
      {every week the team checks each item: nothing enters unchecked}
    \draw[flow] (a) -- (o);
    \draw[flow] (o) -- (l);
    \draw[flow] (l) -- (r);
    % Approved changes
    \node[fbout] (u1) at (13.45,-2.25) {\faQuestionCircle\ \textbf{FAQ entries} $\to$ the index};
    \begin{scope}[on background layer]
      \node[fit=(u1), rounded corners=5pt, draw=good!50, fill=good!3,
            inner sep=5pt] (box) {};
    \end{scope}
    \draw[flow] (r) -- node[right, text=black!70] {approved} (u1);
    \draw[flow] (box.west) -| node[pos=0.25, above, text=black!70] {better answers}
      ([xshift=-0.15cm]a.west) -- (a.west);
  \end{tikzpicture}
  \caption{The operators' corrections, and the questions left without an answer, become better
    answers after a human review.}
  \label{fig:feedback}
\end{figure}

